\section{Check 1: Labels}\label{sec:labels}

\emph{What the benchmark hides.} Each CIC-IoT2023 capture is named after one attack, and all its flows are commonly labeled with that name; but the captures were recorded on a network of active IoT devices, so the attack is only part of the traffic. We call the Benign flows inside an attack capture \emph{background}. Garcia and Valeros predict that models trained with one label per PCAP ``may learn to detect operating system benign traffic as malicious'' (Sec.~3.2)~\cite{garcia2023labeling}; we measure this on traffic the model never saw.

\emph{Evidence labels.} We identified the attackers and victims in every attack capture from the packets themselves (Table~\ref{tab:background}): HTTP payloads against the DVWA (Damn Vulnerable Web Application) test site for the web attacks, destination ports and their timing for brute force, top talkers for the floods, Mirai and the port scan, and ARP replies for MITM-ARP. A flow is labeled as an attack only if (i) one endpoint is a confirmed attacker and the other a victim, on the attack's port where known, or (ii) for ARP spoofing, its first packet was sent from or to the attacker's MAC address and the attacker's own IP is not an endpoint, i.e.\ it is a victim's traffic relayed through the attacker. Every other flow is Benign.

\begin{table}[t]
\centering
\caption{Background share (Benign flows under evidence labels) of each attack capture, and the evidence that confirmed the attacker. Addresses are last octets in 192.168.137.0/24.}
\label{tab:background}
\footnotesize
\setlength{\tabcolsep}{3pt}
\begin{tabular}{lrr>{\raggedright\arraybackslash}p{3.85cm}}
\hline
Capture & Flows & Backgr. & Attacker evidence \\
\hline
XSS          &   4,992 & 98.5\% & \texttt{<script>} from .147 \\
CmdInj       &   5,639 & 98.3\% & DVWA \texttt{exec} page from .147 \\
MITM-ARP     &  12,208 & 97.6\% & MAC \ldots:27:d5 claims .1 and .65 \\
SQLi         &   7,383 & 94.8\% & \texttt{SLEEP(5)}, \texttt{UNION SELECT} \\
Mirai        &     653 & 93.7\% & 7 bots to .191 (UDP) \\
BruteForce   &  11,824 & 90.9\% & .65 to SSH :22, then RTSP :554 \\
Recon        & 205,065 & 30.2\% & .178: 1,001 ports on 58 hosts \\
DoS-HTTP     &  51,275 & 23.0\% & .66 to .71/.141:80 \\
DDoS-SYN     & 511,936 &  0.1\% & 6 sources to .245:4070 \\
\hline
\end{tabular}
\end{table}

\emph{Background share.} Background makes up 90.9--98.5\% of the flows in six of the nine attack captures (web, brute force, MITM-ARP, Mirai) (Table~\ref{tab:background}); overall, file-name and evidence labels differ for 13.4\% of the 857,206 flows.

\emph{The cost of file-name labels on unseen traffic.} We trained the same tree on file-name labels and on evidence labels and measured the Benign FPR. Within a split this is partly circular: the file-label tree is tested on the later background of captures whose earlier background it learned under the capture's name. We therefore also train on nine captures and test on the tenth (leave one capture out, LOCO; Table~\ref{tab:filelabel}). On a held-out attack capture (409--61,921 background flows), the file-label tree flags 71.4--87.8\% of the background flows, against 0.5--5.4\% for the tree trained on evidence labels (the top is the DDoS-SYN capture); the paired difference is 70.9--86.7 points, positive in all nine captures. (Holding out the single Benign capture leaves the file-label tree no Benign training flow: FPR 100\% by construction.) Within the temporal split its pooled FPR of 71.8\% (38,034 of 52,993 Benign test flows) mixes both kinds of traffic: 88.6\% of these false positives are background inside attack captures. Scored against file-name labels, the two trees swap places (0.537 against 0.409), because correct Benign verdicts in attack captures then count as misses. The random split gives the same pattern (pooled FPR 70.2\% against 1.1\%).

\begin{table}[t]
\centering
\caption{The tree trained on file-name labels (file) or on evidence (CIC-IoT2023) or Zeek (IoT-23) labels. Macro-F1 scored against file-name labels (as a file-label study would report) or against evidence/Zeek labels. Benign FPR on background in attack captures and on the benign capture(s). LOCO: range over held-out captures. CIC: CIC-IoT2023. --: not computed. Original labels (corrected values in the text). The two IoT-23 values 0.649 are 0.6486 (file) and 0.6487 (Zeek).}
\label{tab:filelabel}
\footnotesize
\setlength{\tabcolsep}{2.6pt}
\begin{tabular}{lllcccc}
\hline
 & & Trained & \multicolumn{2}{c}{Macro-F1 vs.} & \multicolumn{2}{c}{Benign FPR} \\
Data & Test & on & file & evid./Zeek & backgr. & benign capt. \\
\hline
CIC & temporal & file  & 0.537 & 0.396 & 86.1\% & 31.4\% \\
            & temporal & evid. & 0.409 & 0.725 & 4.8\% & 1.0\% \\
            & LOCO     & file  & --    & --    & 71.4--87.8\% & (100\%) \\
            & LOCO     & evid. & --    & --    & 0.5--5.4\% & 1.5\% \\
\hline
IoT-23      & temporal & file  & 0.649 & 0.169 & 99.97\% & 4.7\% \\
            & temporal & Zeek  & 0.157 & 0.649 & 3.4\% & 0.18\% \\
            & LOCO     & file  & --    & --    & 9.4--99.95\% & 23.4--46.7\% \\
            & LOCO     & Zeek  & --    & --    & 0--49.1\% & 0--17.5\% \\
\hline
\end{tabular}
\end{table}

\emph{Label audit.} Our first DDoS-SYN labeling, from the first 300,000 packets, missed a sixth source (.36) that joins later; error analysis found it as 65,525 ``Benign'' flows predicted as DDoS-SYN, and we re-checked every flood victim over the whole trimmed file. Because that revision moved labels towards the model, we then audited them blind and in both directions: up to 30 test flows from each of 47 strata (capture $\times$ label $\times$ model agreement; 975 flows), rated from their own packets by a fixed program that sees the scenario but not the label, the attacker list or the verdict. Of its first run's 91 disagreements, 67 were rater defects (the RTSP path dictionary of Section~\ref{sec:splits}, flows starting mid-connection, a truncated payload), which we fixed and report. Of the 630 flows the amended rater could decide, 2 of 265 attack labels were wrong (UPnP requests of the SQLi attacker to a non-victim) and 22 of 365 Benign labels, all the DoS-HTTP attacker flooding a third server (.174) missing from our victim list. Weighted by stratum size, the point estimates are 0.004\% of the attack-labeled and 4.3\% of the Benign-labeled test flows (2,269, all that flood); with at most 30 flows per stratum, a stratum without a found error is bounded only below 11--13\% (95\% Wilson). The 345 undecided flows, mostly short or without payload, lie mainly where label and model disagree (239) and weigh 9\% and 14\% of the two groups; errors there cannot be ruled out. MITM-ARP is rated by the MAC relay that defines its label, so for it the audit checks only consistency. The abstract and tables keep the original labels, which understate the evidence-label tree. With both errors corrected, the deployed tree's Benign FPR is 1.3\% instead of 3.8\% (it flagged most of the flood), macro-F1 0.729 [0.585, 0.827], the DoS-HTTP background share 18.6\%, the random--temporal gap $+0.145$ (BruteForce 75\%) and the evidence-label LOCO range 0.7--4.4\%; file-label numbers move by at most 1.3 points.

\emph{Robustness test on IoT-23.} IoT-23 tests the result against labels that are not ours (the authors' Zeek labels). For file-name labels we map each malware family to its main behavior, e.g.\ Hide\&Seek to PortScan, Torii to C\&C and Mirai to DDoS; this mapping is our choice. The background share (Zeek label Benign) of the malware captures ranges from 5.5\% to 99.9\%. The benign-device captures are the out-of-capture test (Table~\ref{tab:filelabel}). Within the temporal split the file-label tree flags 4.7\% of their flows, against 0.18\% for the tree trained on Zeek labels; under LOCO it flags 23.4--46.7\% of a held-out benign device's flows (321--4,577 flows per capture), against 0--17.5\%, more in each of the three (by 7.6--46.7 points). Generalization to an unseen capture is hard regardless of labels: the Zeek-label tree also flags 49.1\% and 42.0\% of the background in the held-out malware captures 3-1 and 42-1. In-capture, the file-label tree flags 97--100\% of the background in every malware capture, a number that mainly shows that it learned each capture's name.

\emph{The benchmark as distributed.} Every study that uses the released CSVs inherits file-name labels. The CSVs (mirror \texttt{bencorn/CIC-IoT-2023}) hold 39 packet-window features without address or timestamp, labeled only by the file name or a \texttt{Label} column equal to it, with rows shuffled, so they cannot be relabeled from evidence. Even scored against file-name labels on a random split, a decision tree and a random forest on these CSVs reach macro-F1 0.749 and 0.803, but BruteForce recall of only 0.43 and 0.41 and Web recall of 0.37 and 0.35, so the weak low-rate classes are not an artifact of our NFStream flows.
